\section{Objectifs}
\begin{itemize}
\item Définir et utiliser une variable aléatoire sur un univers fini.
\item Construire sa loi.
\item Calculer espérance, variance et écart type.
\end{itemize}

\section{Prérequis}
Univers fini, probabilités, sommes pondérées.

\section{Activité de découverte}
Dans un jeu, plusieurs issues peuvent produire le même gain : la variable aléatoire regroupe les issues par valeur numérique.

\section{Vocabulaire}
Variable aléatoire, loi, espérance, variance, écart type.

\section{Cours}
Une variable aléatoire réelle $X$ associe un réel à chaque issue. Sa loi donne les valeurs $x_i$ et les probabilités $p_i=P(X=x_i)$.

\[
E(X)=\sum_i x_ip_i.
\]

La variance et l’écart type sont
\[
V(X)=E(X^2)-E(X)^2,\qquad \sigma(X)=\sqrt{V(X)}.
\]
La linéarité de l’espérance donne $E(aX+b)=aE(X)+b$.

\section{Exemples corrigés}
Si $X$ vaut $0,1,2$ avec probabilités $0,2;0,5;0,3$, alors $E(X)=1,1$.

\section{Méthode}
Construire la loi complète, vérifier que les probabilités somment à 1, puis calculer $E(X)$, $E(X^2)$, $V(X)$ et $\sigma(X)$.

\section{Erreurs fréquentes}
\begin{itemize}
\item Confondre issue et valeur de $X$.
\item Écrire $V(X)=E(X^2)-E(X)$.
\item Croire que l’espérance doit être une valeur possible de $X$.
\end{itemize}

\section{Synthèse}
Une variable aléatoire transforme une expérience aléatoire en grandeur numérique et permet de quantifier moyenne théorique et dispersion.
